% =====================================================================
%  Chapter 3: Pure gauge theory on the lattice
%  Gattringer & Lang, "Quantum Chromodynamics on the Lattice" (LNP 788)
%  Reconstructed from extract/ch03.txt; equations auto-numbered.
% =====================================================================

\chapter{Pure gauge theory on the lattice}

In the last chapter we started the construction of QCD on the lattice. We
discussed a particularly important feature of QCD, namely the self-interaction
of gluons. This self-interaction makes pure gluodynamics, i.e., QCD without
quarks, an interesting, highly nontrivial theory. Pure gluodynamics shows color
confinement, an important property of full QCD. Since gluodynamics is much
easier to handle than QCD with quarks, it is an important subject for studies
of confinement and its underlying mechanisms.

In this chapter we complete the construction of gluodynamics on the lattice by
discussing the integration measure for the link variables. After a section about
gauge invariance and gauge fixing we discuss the potential between two static
quark sources. The static quark potential can then be used to set the scale,
i.e., to determine the value of the lattice spacing $a$ as a function of the
gauge coupling $g$.

Before we come to these topics, let us briefly collect the formulas defining
pure gauge theory on the lattice. Expectation values of observables $O$ (where
$O$ may also be the product of several terms) are given by
\begin{equation}
  \ev{O} = \frac{1}{Z}\int D[U]\,e^{-S_G[U]}\,O[U],
\end{equation}
where the partition function $Z$ is defined as
\begin{equation}
  Z = \int D[U]\,e^{-S_G[U]}.
\end{equation}
The integration measure for the link variables is the product measure
\begin{equation}
  D[U] = \prod_{n\in\Lambda}\prod_{\mu=1}^{4} dU_\mu(n).
\end{equation}
For the gauge field action we introduced the Wilson action
\begin{equation}
  S_G[U] = \frac{\beta}{3}\sum_{n\in\Lambda}\sum_{\mu<\nu}
           \Re\,\tr\!\left[1 - U_{\mu\nu}(n)\right],
\end{equation}
but we again remark that other formulations are being used as well. In (3.4) we
use a widely used abbreviation
\begin{equation}
  \beta = \frac{6}{g^2},
\end{equation}
the so-called inverse coupling. Note that this is the definition of $\beta$ for
the case of SU(3). The expression for the general case of SU($N$) is given in
(3.93).

\section{Haar measure}

In the construction of the QCD lattice action, the link variables $U_\mu(n)$
are introduced as elements of SU(3). We now address the definition of the
individual measures $dU_\mu(n)$ in the product measure of (3.3) which
integrates the link variables $U_\mu(n)$ over the whole group manifold of
SU(3). We consider the problem in a more general way and discuss the measure
$dU$ with $U$ being an element of a Lie group, i.e., a group where the elements
depend continuously on some parameters.

\subsection{Gauge field measure and gauge invariance}

An important restriction for the gauge field measure comes from the gauge
invariance of our theory. In the last chapter we introduced the gauge
transformation of the link variables as
\begin{equation}
  U_\mu(n) \to U'_\mu(n) = \Omega(n)\,U_\mu(n)\,\Omega(n+\hat\mu)^\dagger .
\end{equation}
The group-valued matrices $\Omega(n)$ can be chosen independently at each
lattice site $n$. The action $S_G[U]$ for the gauge field is invariant under
these transformations and we have
\begin{equation}
  S_G[U'] = S_G[U].
\end{equation}
As for any integral, the result of a path integral should be invariant under a
change of variables, in particular under the gauge transformation (3.6). For
the partition function this requirement reads
\begin{equation}
  Z = \int D[U]\,e^{-S_G[U]}
    = \int D[U']\,e^{-S_G[U]}
    = \int D[U']\,e^{-S_G[U']},
\end{equation}
where in the last step we have used the gauge invariance (3.7) of the gauge
action. Comparing the first and the final expression in (3.8) we obtain
\begin{equation}
  D[U] = D[U'].
\end{equation}
Using the fact that $D[U]$ is a product measure (compare (3.3)), we derive the
condition
\begin{equation}
  dU_\mu(n) = dU'_\mu(n) = d\!\left(\Omega(n)\,U_\mu(n)\,\Omega(n+\hat\mu)^\dagger\right),
\end{equation}
for the integration over the individual link variables. This property is one of
the defining properties of the Haar measure.

\subsection{Group integration measure}

The requirement of gauge invariance leads us to a structure which is well-known
in the mathematical literature, the so-called Haar measure. This is a measure
for integration over a continuous compact group $G$. We have derived the
condition (3.10) for the integration over the single group variable
$U_\mu(n)$. Since $\Omega(n)$ and $\Omega(n+\hat\mu)$ can be chosen
independently, the measure $dU$ for a group element must be invariant under
left- and right multiplication with another group element $V \in G$, i.e.,
\begin{equation}
  dU = d(UV) = d(VU).
\end{equation}
The last equation, together with the normalization (integration over all of
$G$)
\begin{equation}
  \int dU\,1 = 1,
\end{equation}
are the defining properties of the Haar measure for the integration over
compact Lie groups. As announced, we discuss the measure not only for SU(3),
but for more general compact Lie groups $G$, such as SU($N$) or U(1). For these
cases we now give an explicit construction of the Haar measure. Let
$U = U(\omega)$ be an element of the compact Lie group $G$ in the exponential
representation (A.2), parameterized by the set of real numbers
$\omega^{(k)}$. We have already discussed (see also Appendix A.1) that
$\partial/\partial\omega^{(k)}\, U(\omega)\, U(\omega)^{-1}$ is in the Lie
algebra of the group. Based on these Lie algebra-valued objects we can define a
metric $ds^2$ on the group as (we use summation convention for $n$, $m$)
\begin{equation}
  \begin{split}
    ds^2 &\equiv \tr\!\left[
      \frac{\partial U(\omega)}{\partial \omega^{(n)}} U(\omega)^{-1}
      \left(\frac{\partial U(\omega)}{\partial \omega^{(m)}} U(\omega)^{-1}\right)^{\!\dagger}
    \right] d\omega^{(n)} d\omega^{(m)}\\
    &= g(\omega)_{nm}\, d\omega^{(n)} d\omega^{(m)} .
  \end{split}
\end{equation}
From this equation we read off a metric tensor $g(\omega)$,
\begin{equation}
  g(\omega)_{nm} = \tr\!\left[
    \frac{\partial U(\omega)}{\partial \omega^{(n)}} U(\omega)^{-1}
    \left(\frac{\partial U(\omega)}{\partial \omega^{(m)}} U(\omega)^{-1}\right)^{\!\dagger}
  \right]
  = \tr\!\left[
    \frac{\partial U(\omega)}{\partial \omega^{(n)}}
    \frac{\partial U(\omega)^{\dagger}}{\partial \omega^{(m)}}
  \right].
\end{equation}
In terms of the metric $g(\omega)$ we can define the measure $dU$ as
\begin{equation}
  dU = c\,\sqrt{\det[g(\omega)]}\prod_k d\omega^{(k)}.
\end{equation}
Because all groups we consider are compact groups, the parameters
$\omega^{(k)}$ need only be integrated over finite intervals for $U(\omega)$ to
cover all group elements. The constant $c$ in (3.15) can be used to fulfil the
normalization condition (3.12). What is left to show is that the measure
defined by (3.14) and (3.15) is indeed invariant under transformations within
the group, i.e., that it obeys (3.11). Let $U(\tilde\omega)$ be the group
element that is obtained from $U(\omega)$ by left- or right multiplication with
some other group element $V$. The parameters $\tilde\omega^{(k)}$ are related
to the original parameters $\omega^{(k)}$ by some functional relation, i.e., we
can write $\omega^{(k)} = \omega^{(k)}(\tilde\omega)$. The product measure for
the $\omega^{(k)}$ transforms as
\begin{equation}
  \prod_k d\omega^{(k)} = \det[J]\prod_k d\tilde\omega^{(k)},
\end{equation}
where the Jacobi determinant is given by
\begin{equation}
  J_{kn} = \frac{\partial \omega^{(k)}}{\partial \tilde\omega^{(n)}}.
\end{equation}
The metric tensor $g(\tilde\omega)$ for $U(\tilde\omega)$ is given by
\begin{equation}
\begin{split}
  g(\tilde\omega)_{nm} &= \tr\!\left[
    \frac{\partial U(\tilde\omega)}{\partial \tilde\omega^{(n)}}
    \frac{\partial U(\tilde\omega)^{\dagger}}{\partial \tilde\omega^{(m)}}
  \right] \\
  &= \tr\!\left[
    \frac{\partial U(\omega)}{\partial \omega^{(k)}}
    \frac{\partial U(\omega)^{\dagger}}{\partial \omega^{(l)}}
  \right]\frac{\partial \omega^{(k)}}{\partial \tilde\omega^{(n)}}
   \frac{\partial \omega^{(l)}}{\partial \tilde\omega^{(m)}}
  = g(\omega)_{kl}\,J_{kn} J_{lm}.
\end{split}
\end{equation}
In matrix notation the last equation reads ($T$ denotes transposition)
\begin{equation}
  g(\tilde\omega) = J^T g(\omega)\,J .
\end{equation}
Using the factorization property of determinants we obtain from this equation
\begin{equation}
  \det[g(\tilde\omega)] = \det[g(\omega)]\,\det[J]^2 .
\end{equation}
Putting things together, we find that our group measure is indeed invariant,
\begin{equation}
\begin{split}
  d\tilde U &= c\,\sqrt{\det[g(\tilde\omega)]}\prod_k d\tilde\omega^{(k)}
  = c\,\sqrt{\det[g(\omega)]}\,\det[J]\prod_k d\tilde\omega^{(k)} \\
  &= c\,\sqrt{\det[g(\omega)]}\prod_k d\omega^{(k)} = dU .
\end{split}
\end{equation}
In the first step we used (3.20) and in the second step (3.16). This completes
the proof of invariance of the Haar measure.

\subsection{A few integrals for SU(3)}

In order to become familiar with the properties of the Haar measure, we now
discuss a few SU(3) integrals which we will need subsequently. Our discussion
is based on the invariance of the measure $dU$. The integrals we study are
integrals over products of entries $U_{ab}$ of group elements $U$ in the
fundamental representation. In particular, we consider the integrals
\begin{equation}
  \int_{SU(3)} dU\; U_{ab} = 0 ,
\end{equation}
\begin{equation}
  \int_{SU(3)} dU\; U_{ab}\, U_{cd} = 0 ,
\end{equation}
\begin{equation}
  \int_{SU(3)} dU\; U_{ab}\, (U^\dagger)_{cd} = \frac{1}{3}\,\delta_{ad}\,\delta_{bc} ,
\end{equation}
\begin{equation}
  \int_{SU(3)} dU\; U_{ab}\, U_{cd}\, U_{ef} = \frac{1}{6}\,\varepsilon_{ace}\,\varepsilon_{bdf} .
\end{equation}
The basic tool for analyzing these integrals is the following equation for
integrals over functions $f(U)$:
\begin{equation}
  \int_{SU(3)} dU\; f(U)
  = \int_{SU(3)} dU\; f(VU)
  = \int_{SU(3)} dU\; f(UW) .
\end{equation}
The elements $V$ and $W$ are arbitrary SU(3) matrices. This relation follows
directly from the invariance (3.11) of the measure. Applying (3.26) to the
left-hand side of the first integral (3.22) we obtain
\begin{equation}
  \int_{SU(3)} dU\; U_{ab}
  = \int_{SU(3)} dU\; (VU)_{ab}
  = V_{ac}\int_{SU(3)} dU\; U_{cb} ,
\end{equation}
where the index $c$ is summed over. In order to fulfil this equation in a
nontrivial way, we must have $V_{ac} = \delta_{ac}$. However, the last equation
must hold for arbitrary group elements $V$, such that the integral itself must
vanish, thus establishing (3.22). In an equivalent way one shows (3.23).

The Eq.~(3.24) is somewhat more interesting. Let us for the moment set $b = c$
in the integrand and sum over $b$. We obtain
\begin{equation}
  \begin{split}
    \int_{SU(3)} dU\;\big[ U_{a1}(U^\dagger)_{1d} + U_{a2}(U^\dagger)_{2d} + U_{a3}(U^\dagger)_{3d} \big]
    &= \int_{SU(3)} dU\; (UU^\dagger)_{ad}\\
    &= \int_{SU(3)} dU\; \delta_{ad} = \delta_{ad} .
  \end{split}
\end{equation}
In the second step we have used $U^\dagger = U^{-1}$ and in the last step the
normalization of the Haar measure (3.12). For the left-hand side of (3.28) it
is important to realize that the three contributions are entirely equivalent.
For example the first term can be transformed into the second term by
exchanging rows 1 and 2 as well as columns 1 and 2. This operation transforms
an element of SU(3) into another group element, and both of these elements are
summed over in the group integral. This equivalence implies that each of the
three terms on the left-hand side contributes one-third to the result on the
right-hand side of (3.28). Thus we have established formula (3.24) for the case
$b = c$ (not summed). For $b \neq c$ we can again apply (3.26) to show that for
this case the integral has to vanish. The essence of (3.26) is that the
integral only gives a nonvanishing result if the integrand contributes to 1,
i.e., the Haar measure projects out the contribution of the integrand to the
so-called trivial or singlet representation of the group.

From (3.24) immediately follows a relation that will be useful later:
\begin{equation}
  \int dU\; \tr[VU]\,\tr[U^\dagger W] = \frac{1}{3}\,\tr[VW] .
\end{equation}
This equation allows one to integrate the common link variable that occurs in a
product of two plaquettes, resulting in a trace of link variables around the
two plaquettes. This result is depicted in Fig.~3.1.

\begin{figure}[htbp]
  \centering
  \includegraphics[width=0.7\textwidth]{fig31.png}
  \caption*{Fig. 3.1. \emph{Integrating out the common link of a product of two
  plaquettes}.}
  \label{fig:31}
\end{figure}

Essentially the same mechanism as in (3.24) is also at work in the integral
(3.25): Again, different possible combinations of the integrand can be summed
to produce a singlet. In particular, if one sets the second indices $b, d, f$
of the integrand to the values $b = 1$, $d = 2$, $f = 3$, and sums the first
indices $a, c, e$ with the completely anti-symmetric $\varepsilon$-tensor, one
finds
\begin{equation}
  \varepsilon_{ace}\, U_{a1}\, U_{c2}\, U_{e3} = \det[U] = 1 ,
\end{equation}
where we have used the definition of the determinant and the fact that elements
of SU(3) have $\det[U] = 1$. If two of the second indices in the integrand of
(3.30) are equal, say $b = d = 1$, then summing with the $\varepsilon$-tensor
corresponds to a determinant with two equal columns, and therefore this sum
vanishes. If two of the second indices are interchanged with respect to (3.30),
e.g., $b = 2$, $d = 1$ then this corresponds to the interchange of two rows and
the determinant acquires an extra minus sign. Thus, we obtain the following
formula for summing the integrand of (3.30) with the $\varepsilon$-tensor
\begin{equation}
  \varepsilon_{ace}\, U_{ab}\, U_{cd}\, U_{ef} = \varepsilon_{bdf} .
\end{equation}
Thus, whenever the set of the first coefficients $a, c, e$ and the set $b, d,
f$ both give rise to nonvanishing $\varepsilon$-tensor contributions
($\varepsilon_{ace} \neq 0$, $\varepsilon_{bdf} \neq 0$), the integrand gives a
contribution to the singlet. The sign of this contribution is given by the
combination of $\varepsilon$-tensors on the right-hand side of (3.25). If this
product of $\varepsilon$-tensors vanishes for the given combination $a, b, c,
d, e, f$, the integrand does not contribute to the singlet and the integral
vanishes. Using the same argument as before, based on the simultaneous
interchange of rows and columns, one finds that all the nonvanishing
contributions have the same weight. The fact that there are six terms on the
left-hand side of (3.31) then determines the factor $1/6$ on the right-hand
side of (3.25).

This completes our excursion into group integrals. The integrals (3.22),
(3.23), (3.24), (3.25), and (3.29) can easily be generalized to the case of
SU($N$). An introductory account of more general SU($N$) group integrals can be
found in the book by Creutz [1]. We finally remark that for a treatment of the
strong coupling expansion in more detail than in our Sect.~3.4, the concept of
character expansion is a powerful tool. However, character expansion, as well
as the general representation theory of SU($N$), is beyond the scope of this
book and we refer the reader to [2--5].

\section{Gauge invariance and gauge fixing}

We have seen that gauge invariance is one of the central principles in the
construction of QCD. In this section we discuss the freedom of choosing a gauge
in more detail and address the different roles of gauge fixing in the continuum
and on the lattice. We argue that physical observables have to be
gauge-invariant. We stress that all the gauge properties discussed here for
pure gauge theory can be taken over to the full theory with fermions.

\subsection{Maximal trees}

Let us begin our discussion of gauge fixing with a kind of gauge that is
particular for the lattice. We show that in a certain set of link variables, a
so-called maximal tree, the link variables can be set to the identity element
1 and left out in the integration of the gauge field.

We start with an arbitrary configuration of link variables $U_\mu(n)$. In the
beginning of our construction all gauge transformation matrices $\Omega(n)$ are
set to the identity 1. We pick a single-link variable $U_{\mu_0}(n_0)$ and set
the transformation matrix $\Omega(n_0 + \hat\mu_0)$ at the endpoint of the link
to the value $U_{\mu_0}(n_0)$, keeping all other $\Omega(n)$ at 1. Our link
$U_{\mu_0}(n_0)$ thus is transformed to
\begin{equation}
  U'_{\mu_0}(n_0)
  = \Omega(n_0)\,U_{\mu_0}(n_0)\,\Omega(n_0+\hat\mu_0)^\dagger
  = 1\,U_{\mu_0}(n_0)\,U_{\mu_0}(n_0)^\dagger = 1 .
\end{equation}
The transformation with the nontrivial $\Omega(n_0+\hat\mu_0)$ will also affect
all link variables starting at the site $n_0+\hat\mu_0$, in particular
\begin{equation}
  U'_{\mu_1}(n_0+\hat\mu_0) = \Omega(n_0+\hat\mu_0)\,U_{\mu_1}(n_0+\hat\mu_0) .
\end{equation}
For these transformed links we can repeat the step of (3.32) and choose the
matrices $\Omega$ at their endpoints such that also these links are transformed
to the identity 1. The whole procedure can be repeated until we hit a link
$U_{\mu'}(n')$ which connects to another link that already has been transformed
to 1 before. If we wanted to transform also this particular link
$U_{\mu'}(n')$, according to (3.33) we would transform the other link, which is
already at 1, away from the identity. This restricts the set of links that can
be transformed to 1 to a cluster of links which does not contain closed loops.
This requirement defines a maximal tree as a maximal collection of link
variables that can be transformed to 1 (note that there are many different
maximal trees). In Fig.~3.2 we show an example of a maximal tree on a 2D
sublattice. Note, however, that also smaller subsets than maximal trees can be
gauged to 1.

\begin{figure}[htbp]
  \centering
  \includegraphics[width=0.6\textwidth]{fig32.png}
  \caption*{Fig. 3.2. \emph{Sketch of a maximal tree on a 2D sublattice. The
  fat lines represent link variables that are set to 1 and omitted in the
  integration}.}
  \label{fig:32}
\end{figure}

Let us now discuss the consequences of the freedom to fix the gauge on a
maximal tree or a subset of a maximal tree. We consider the vacuum expectation
value of some gauge-invariant observable $O$ as defined in (3.1), (3.2), and
(3.3). Gauge invariance of $O$ implies
\begin{equation}
  O[U] = O[U'] ,
\end{equation}
where the link variables in the configurations $U$ and $U'$ are related by a
gauge transformation (3.6). We will discuss the need for gauge invariance of
our observables in more detail below. Also the action and the measure are
gauge-invariant, i.e.,
\begin{equation}
  S_G[U] = S_G[U'] , \qquad D[U] = D[U'] .
\end{equation}
Equations (3.34) and (3.35) imply that the whole integrand and the measure in
(3.1) are unchanged when setting the links in a maximal tree to 1. Since the
construction of a maximal tree works for any particular choice of link
variables, we can keep the links in the maximal tree at 1 throughout the whole
integration $D[U]$. Since the Haar measure is normalized to 1 (see (3.12)), the
links in the maximal tree can be omitted in the integration altogether. Thus,
we can summarize the procedure of fixing the gauge to a maximal tree as
follows: Select a maximal tree, or a smaller set of links without closed loops,
and set the links in this set to 1. Subsequently, integrate over all link
variables that are not contained in the chosen set. The expectation value of a
gauge-invariant observable is unaffected, whether we fix the gauge or not.

The role of gauge fixing on the lattice is very different from the role it
plays in the continuum. Fixing the gauge on the lattice is a step which we can
implement to simplify some calculations or to make the interpretation of
observables more transparent (see the discussion of the Wilson loop below).
However, fixing the gauge is not a necessary step to make vacuum expectation
values of operators well defined and computable. In the continuum the situation
is different: There the kernel of the quadratic part of the gluon action (the
inverse gluon propagator) has zero modes caused by pure gauges, i.e., gauge
transformations of the trivial gauge field configuration. These modes give rise
to singularities in the gluon propagator that have to be removed by choosing a
particular gauge. Only after that step can the continuum theory be well defined
in perturbation theory.

We have already mentioned that the procedure of fixing the gauge via a maximal
tree is special to the lattice formulation. There is, however, a certain gauge,
the so-called temporal gauge, which has a corresponding counterpart in the
continuum. On a lattice with infinite temporal extent we can set all time-like
links to 1, i.e., we can set
\begin{equation}
  U_4(n) = 1 \quad \forall n .
\end{equation}
We remark that this is not a maximal tree but a smaller set. In the (Euclidean)
continuum theory the temporal gauge is defined by $A_4(x) = 0$, which matches
the lattice definition due to the relation $U_4 = \exp(\im a A_4)$. The temporal
gauge is important for the Hamiltonian formulation of continuum gauge theories
(see any textbook on quantum field theory, such as [6--8]). We will make use of
the temporal gauge when we introduce and discuss the Wilson loop.

\subsection{Other gauges}

Other gauges that are used in the continuum also play a role on the lattice, in
particular when one matches lattice results to continuum calculations done in a
particular gauge (e.g., Landau gauge, Coulomb gauge). Also in lattice studies
of the confinement mechanism special gauges, such as the maximally abelian
gauge, are used. These gauges are not implemented via maximal trees but
typically through extremizing some functional. As an example, which illustrates
how such gauges are constructed, we briefly discuss the Landau gauge, which we
will need in Chap.~11 when we discuss renormalization.

In the continuum the Landau gauge is defined by the condition
\begin{equation}
  \partial_\mu A_\mu(x) = 0 .
\end{equation}
The important step is to translate this prescription, which in the continuum is
formulated in terms of the algebra-valued gauge field $A_\mu(x)$, into a
prescription that can be used on the lattice, where we work with the
group-valued link variables $U_\mu(n)$. In the first step we show that
requiring the Landau gauge condition is equivalent to finding an extremal value
of the functional
\begin{equation}
  W[A] = \sum_{\mu=1}^{4}\int \dd^4 x\, \tr\!\left[A_\mu(x)^2\right] ,
\end{equation}
with regard to gauge transformations. If $W[A]$ is at an extreme value, it has
to be invariant when we apply an infinitesimal gauge transformation defined by
$\Omega(x) = \exp(\im\,\varepsilon\, H(x))$. Note that $H(x)$ is a traceless
hermitian matrix, such that $\Omega(x)^\dagger = \exp(-\im\,\varepsilon\,
H(x))$. A line of algebra shows that, for our infinitesimal transformation, the
gauge transformation (2.11) reduces to
\begin{equation}
  A_\mu(x) \to A_\mu(x) + \varepsilon\,\big(\im\,[H(x), A_\mu(x)] - \partial_\mu H(x)\big) + \order{\varepsilon^2} .
\end{equation}
Inserting the transformed gauge field into the functional $W[A]$, one finds
that it changes according to
\begin{equation}
\begin{split}
  W[A] &\to W[A] - 2\varepsilon\sum_{\mu=1}^{4}\int \dd^4 x\,
    \tr\!\left[A_\mu(x)\,\partial_\mu H(x)\right] + \order{\varepsilon^2} \\
  &= W[A] + 2\varepsilon\sum_{\mu=1}^{4}\int \dd^4 x\,
    \tr\!\left[(\partial_\mu A_\mu(x))\, H(x)\right] + \order{\varepsilon^2} ,
\end{split}
\end{equation}
where the periodicity of the trace makes the contribution of the commutator in
(3.39) vanish. In the second step we performed an integration by parts
(assuming that $H(x)$ vanishes for large $x$). If $W[A]$ is at an extremum, then
the $\order{\varepsilon}$ term in (3.40) has to vanish. Since $H(x)$ can be
chosen as an arbitrary traceless hermitian matrix, this statement is equivalent
to (3.37).

Transforming the $W[A]$ into a corresponding lattice prescription is simple. We
can consider the functional
\begin{equation}
  W_{\mathrm{lat}}[U] = -a^2\sum_{n\in\Lambda}\sum_{\mu=1}^{4}
    \tr\!\left[U_\mu(n) + U_\mu(n)^\dagger\right] ,
\end{equation}
which upon expansion of $U_\mu(n) = \exp(\im a A_\mu(n))$ in small $a$ turns
into $W[A]$ (the additional constant is irrelevant).

The goal is to find a gauge transformation $\Omega$ which shifts
$W_{\mathrm{lat}}[U]$ to a minimal value for a given gauge configuration. More
explicitly we have to find the minimum of
\begin{equation}
  F[\Omega] = -a^2\sum_{n\in\Lambda}\sum_{\mu=1}^{4}
    \tr\!\left[\Omega(n)U_\mu(n)\Omega(n+\hat\mu)^\dagger
    + \Omega(n+\hat\mu)\,U_\mu(n)^\dagger\,\Omega(n)^\dagger\right] ,
\end{equation}
as a function of the matrices $\Omega(n)$ for fixed $U_\mu(n)$. This is a
large-scale optimization problem which can be attacked with methods such as
overrelaxation or simulated annealing (see, e.g., [9]). A particular problem is
the constraint that the matrices $\Omega(n)$ have to be in SU(3). We discuss
how such matrices are generated numerically in Chap.~4. We remark that a
straightforward modification (sum $\mu$ from 1 to 3 only) of the method
presented here for Landau gauge allows one to implement the Coulomb gauge
condition $\sum_{j=1}^{3}\partial_j A_j = 0$ on the lattice. For a review of
gauge-fixing procedures on the lattice see [10].

\subsection{Gauge invariance of observables}

In our discussion of maximal trees we already assumed that physical observables
$O$ are gauge-invariant, i.e., that they obey (3.34). We now discuss this
requirement in more detail.

By inspecting the integral (3.22) one sees that the attempt to use an individual
link variable as an observable fails because its expectation value vanishes.
Such an individual link variable is an example of a non-gauge-invariant
functional $O[U]$, i.e., it does not obey $O[U] = O[U']$ for configurations $U$
and $U'$ related by a gauge transformation. The expectation value of any
non-gauge-invariant functional is equal to the expectation of its average over
the gauge group. To see this we perform a gauge transformation
$U_\mu(n) \to \Omega(n)U_\mu(n)\Omega(n+\hat\mu)^\dagger$ and obtain an
expectation value of an arbitrary observable
\begin{equation}
  \begin{split}
    \ev{O} &= \frac{1}{Z}\int D[U]\,e^{-S_G[U]}\,O[U]\\
    &= \frac{1}{Z}\int D[\Omega U\Omega^\dagger]\,e^{-S_G[\Omega U\Omega^\dagger]}
      \,O[\Omega U\Omega^\dagger]\\
    &= \frac{1}{Z}\int D[U]\,e^{-S_G[U]}\,O[\Omega U\Omega^\dagger] ,
  \end{split}
\end{equation}
where in the last step we used both the gauge invariance of the action and the
invariance of the Haar measure. Since the values $\Omega(n)$ can be chosen
arbitrarily, we may even average over the gauge group:
\begin{equation}
  \ev{O} = \frac{1}{Z}\int D[U]\,e^{-S_G[U]}
    \int D[\Omega]\,O[\Omega U\Omega^\dagger]
  = \ev{\int D[\Omega]\,O[\Omega U\Omega^\dagger]} .
\end{equation}
Following these lines of argument, Elitzur [11] proved that local gauge
symmetries cannot be broken spontaneously.

According to (3.44) we find for our example of a single-link variable
\begin{equation}
  \ev{U_\mu(n)}
  = \int d\Omega(n)\,\ev{\Omega(n)U_\mu(n)}
  = \int d\Omega(n)\,\Omega(n)\,\ev{U_\mu(n)} = 0 .
\end{equation}
We remark that for some calculations it is indeed advantageous to work with
observables that are not gauge-invariant and to evaluate them after the gauge
has been fixed. As mentioned, this procedure is useful when lattice results are
compared to a continuum calculation done in a particular gauge.

\section{Wilson and Polyakov loops}

In this section we present observables which allow one to determine the
potential between two static color sources. These observables are the so-called
Wilson and Polyakov loops which we first introduce and only later give their
interpretation.

\subsection{Definition of the Wilson loop}

In the last section we discussed that physical observables have to be
gauge-invariant. A prototype of a gauge-invariant object, made from only the
gauge fields, is the trace of a product of link variables along a closed loop
which we introduced in (2.46),
\begin{equation}
  L[U] = \tr\!\left[\prod_{(n,\mu)\in L} U_\mu(n)\right] .
\end{equation}
Here $L$ is a closed loop of links on the lattice and the product in (3.46)
runs over all these links. The Wilson loop which we introduce now is of that
type. A Wilson loop $W_L$ is made from four pieces, two so-called Wilson lines
$S(m, n, n_t)$, $S(m, n, 0)$, and two temporal transporters $T(n, n_t)$,
$T(m, n_t)$. The Wilson line $S(m, n, n_t)$ connects the two spatial points $m$
and $n$ along some path $C_{m,n}$ with all link variables restricted to time
argument $n_t$,
\begin{equation}
  S(m, n, n_t) = \prod_{(k,j)\in C_{m,n}} U_j(k, n_t) .
\end{equation}
The temporal transporter $T(n, n_t)$ is a straight line of $n_t$ link variables
in time direction, all situated at spatial position $n$,
\begin{equation}
  T(n, n_t) = \prod_{j=0}^{n_t-1} U_4(n, j) .
\end{equation}
Attaching the four pieces to each other gives a closed loop $L$,
\begin{equation}
  L: \quad (m, n_t) \xrightarrow{S} (n, n_t) \xrightarrow{T^\dagger} (n, 0)
     \xrightarrow{S^\dagger} (m, 0) \xrightarrow{T} (m, n_t) .
\end{equation}
The Wilson loop $W_L$ is obtained by taking the trace,
\begin{equation}
  W_L[U] = \tr\!\left[S(m, n, n_t)\,T(n, n_t)^\dagger\,S(m, n, 0)^\dagger\,
    T(m, n_t)\right]
  = \tr\!\left[\prod_{(k,\mu)\in L} U_\mu(k)\right] .
\end{equation}
If the piece of loop $C_{m,n}$ used in $S(m, n, n_t)$ is a straight line we
speak of a planar Wilson loop. Note that this can be the case only if $m$ and
$n$ fall on a common coordinate axis. Otherwise the Wilson loop is called
nonplanar. Figure 3.3 shows an example of a planar and a nonplanar loop.

\begin{figure}[htbp]
  \centering
  \includegraphics[width=0.8\textwidth]{fig33.png}
  \caption*{Fig. 3.3. \emph{Examples for a planar (left-hand side plot) and a
  nonplanar (right-hand side) Wilson loop. The horizontal direction is time}.}
  \label{fig:33}
\end{figure}

\subsection{Temporal gauge}

Before we can discuss the physical interpretation of the Wilson loop we must
discuss a peculiarity of gauge theories. If one wants to evaluate the canonical
momentum (1.45) for the gauge field action (2.17) one finds that for the
temporal component $A_4$ the canonical momentum vanishes. The reason is that
the field strength tensor $F_{\mu\nu}(x) = -\im\,[D_\mu(x), D_\nu(x)]$ does not
contain derivatives of $A_4$ with respect to time. A possible way out of this
problem is to use a gauge where
\begin{equation}
  A_4(x) = 0 ,
\end{equation}
i.e., the aforementioned temporal gauge. We remark that simply stating (3.51)
does not by far do justice to the subtleties involved in the quantization of
gauge theories, and we refer the reader to field theory books such as [6--8] on
this issue. On the lattice, temporal gauge corresponds to the condition (3.36).

We stress that in the following we use the temporal gauge only to find the
physical interpretation of the Wilson loop. For the actual computation of the
expectation value we do not need to fix the gauge. The result for the
expectation value of the Wilson loop is of course the same whether we fix the
gauge or not.

\subsection{Physical interpretation of the Wilson loop}

In the temporal gauge (3.36), discussed in the last paragraph, the temporal
transporters become trivial,
\begin{equation}
  T(n, n_t) = \prod_{j=0}^{n_t-1} U_4(n, j) = 1 ,
\end{equation}
and we obtain the following chain of identities
\begin{equation}
  \ev{W_L} = \ev{W_L}_{\mathrm{temp}}
  = \ev{\tr\!\left[S(m, n, n_t)\,S(m, n, 0)^\dagger\right]}_{\mathrm{temp}} ,
\end{equation}
where in the first step we have used the fact that the expectation value of a
gauge-invariant observable remains unchanged when fixing the gauge. In the
second step we used $W_L[U] = \tr[S(m, n, n_t)\,S(m, n, 0)^\dagger]$, which
follows from the definition of the Wilson loop in (3.50) and the triviality of
the temporal transporter in temporal gauge (3.52).

The temporal gauge used in (3.53) makes explicit that the Wilson loop is the
correlator of two Wilson lines $S(m, n, n_t)$ and $S(m, n, 0)$ situated at time
slices $n_t$ and 0. Thus we can interpret this correlator using our first key
equation in the form of (1.21). Accordingly, the correlator behaves for large
total temporal extent $T$ of our Euclidean lattice as ($a$, $b$ are summed)
\begin{equation}
  \ev{\tr\!\left[S(m, n, n_t)\,S(m, n, 0)^\dagger\right]}_{\mathrm{temp}}
  = \sum_k \bra{0}\tilde S(m,n)_{ab}\ket{k}\,
    \bra{k}\tilde S(m,n)^\dagger_{ba}\ket{0}\, e^{-tE_k} ,
\end{equation}
where the Euclidean time argument $t$ is related to $n_t$ via $t = a\,n_t$ with
$a$ being the lattice spacing. The sum in (3.54) runs over all states $\ket{k}$
that have a nonvanishing overlap with $\tilde S(m,n)^\dagger\ket{0}$.

In the next paragraph we will argue that the states $\ket{k}$ with nonvanishing
overlap are states describing a static quark--antiquark pair located at spatial
positions $m$ and $n$. Thus in (3.54) the term with the lowest energy $E_1$ is
expected to be the state describing our static quark--antiquark pair. Higher
states could be, e.g., this pair plus additional particle--antiparticle
combinations with the quantum numbers of the vacuum. The energy $E_1$ is thus
identified with the energy of the quark--antiquark pair, which is the static
potential $V(r)$ at spatial quark separation $r$,
\begin{equation}
  E_1 = V(r) \qquad \text{with} \quad r = a\,|m - n| .
\end{equation}
Combining (3.53), (3.54), and (3.55) we obtain
\begin{equation}
  \ev{W_L} \propto e^{-tV(r)}\left(1 + \order{e^{-t\Delta E}}\right)
  = e^{-n_t a V(r)}\left(1 + \order{e^{-n_t a \Delta E}}\right) .
\end{equation}
Thus we find that we can calculate the static quark--antiquark potential from
the large-$n_t$ behavior of the Wilson loop. The corrections in (3.56) are
exponentially suppressed, where $\Delta E$ is the difference between $V(r)$ and
the first excited energy level of the quark--antiquark pair.

We stress that Wilson loops are oriented. However, reversing the orientation,
which on an algebraic level is complex conjugation (compare (2.34)), simply
corresponds to the interchange of quark and antiquark. Thus both orientations
serve equally well for a determination of the potential $V(r)$.

The Wilson loops we have introduced are not necessarily planar. In Fig.~3.3 we
show two Wilson loops, a planar one (left-hand side plot) and a nonplanar loop
(right-hand side). Both loops have $n_t = 5$ (the horizontal direction is
time). The planar loop has $r = 3a$, the nonplanar loop has
$r = \sqrt{3^2 + 1}\,a = \sqrt{10}\,a$. Thus with nonplanar Wilson loops we can
calculate the potential $V(r)$ not only at distances $r$ that are integer
multiples of $a$, but also at intermediate points. Nonplanar Wilson loops also
allow one to study whether rotational invariance is eventually restored when
approaching the continuum limit.

Wilson loops may also be used as operators for purely gluonic bound states, the
so-called glueballs. The exponential decay of the corresponding Euclidean
correlation functions allow to extract the mass spectrum of these states (cf.
Chap.~6 and the discussion in [12]).

\subsection{Wilson line and the quark--antiquark pair}

To complete our physical interpretation of the Wilson loop we still need to show
that the states $\tilde S(m,n)^\dagger_{ba}\ket{0}$ do indeed have overlap with
a quark--antiquark pair. However, the true derivation that $S(m, n)$ is the
correct expression to describe a quark--antiquark pair in the limit of large
quark mass, will only be given in Chap.~5 after we have provided the final
definition of the quark fields on the lattice. There we will show that the quark
propagator reduces to $S(m, n)$ in the limit of infinitely heavy quarks. For
now we are content with showing that $S(m, n)$ has the same transformation
properties as a quark--antiquark pair under a gauge transformation.

According to our discussion in Sect.~2.1, a quark--antiquark pair at spatial
positions $am$, $an$ is described by a product of fields
\begin{equation}
  Q(m, n)_{ab} \equiv \psi(m)^{\alpha}_{a}\,\psi(n)^{\beta}_{b} .
\end{equation}
The quark fields carry spinor ($\alpha$, $\beta$) and color ($a$, $b$) indices.
However, here we are not interested in the dependence of the potential on the
spinor indices, and we ignore them in the definition of $Q(m, n)_{ab}$.
$Q(m, n)_{ab}$ is not gauge-invariant. According to (2.7), it transforms under
gauge transformations as
\begin{equation}
  Q(m, n)_{ab} \to \Omega(m)_{aa'}\,Q(m, n)_{a'b'}\,\Omega(n)^\dagger_{b'b} .
\end{equation}
From the discussion following the gauge transformation properties of link
variables (2.33), we know that the products $P[U]$ of link variables, defined
in (2.43), transform exactly as required (compare (2.44)),
\begin{equation}
  S(m, n)_{ab} \to \Omega(m)_{aa'}\,S(m, n)_{a'b'}\,\Omega(n)^\dagger_{b'b} .
\end{equation}
Thus we have at least verified that the Wilson line has the same transformation
properties as the quark--antiquark pair.

\subsection{Polyakov loop}

Let us conclude this section with discussing a modification of the Wilson loop,
the so-called Polyakov loop [13] (also called thermal Wilson line). Here we
work with boundary conditions for the gauge fields that are periodic in the
time direction. We make the temporal extent $n_t$ of the Wilson loop as large as
possible on our lattice, i.e., we set $n_t = N_T$, where $N_T$ is the total
number of lattice points in time direction. Then the spatial pieces of the
Wilson loop (the non-horizontal lines in Fig.~3.3) sit on top of each other but
are oriented in opposite direction. Due to the periodic boundary conditions we
cannot gauge-transform all temporal links to 1. We can, however, gauge the
spatial pieces of our loop to 1. Then the Wilson loop reduces to the two
disconnected paths (compare (3.48)) $T(m, N_T)$, $T(n, N_T)^\dagger$ of temporal
link variables, located in space at the two positions $m$ and $n$. Both these
paths wind around the temporal direction of the lattice but have opposite
orientations.

We can make this new observable gauge-invariant by taking the trace for each of
the two loops individually. This is simply a rearrangement of the color indices
and leaves the interpretation of the observable the same. In this way we
introduce the so-called Polyakov loop
\begin{equation}
  P(m) = \tr\!\left[\prod_{j=0}^{N_T-1} U_4(m, j)\right] ,
\end{equation}
which as a trace over a closed loop is gauge-invariant. We now can abandon our
special gauge and obtain ($r = a|m - n|$)
\begin{equation}
  \ev{P(m)P(n)^\dagger} \propto e^{-N_T a V(r)}\left(1 + \order{e^{-N_T a\Delta E}}\right) .
\end{equation}
The numerical calculation of the static potential which we present in Chap.~4
is based on the Polyakov loop, i.e., on the two Eqs.~(3.60) and (3.61). The
identity of the potential as defined via Wilson or Polyakov loops,
respectively, is not rigorously proven; there is a proof that the string tension
derived from the Polyakov loop correlator is bounded from above by that of the
Wilson loop [14].

An alternative way of motivating the introduction of the Polyakov loop is to
couple a current $j_\mu$ to the field $A_\mu$ by defining an operator
$O = \tr\,P\exp\!\big(\im\int \dd^4 z\, j_\mu(z)\, A_\mu(z)\big)$. Using the
current corresponding to a static charge at position $x$, given by
$j_\mu(z) = (0, 0, 0, 1)\,\delta(z - x)$ (Euclidean metric), on the lattice the
operator $O$ translates to our Polyakov loop.

We finally remark that the vacuum expectation value $\ev{P(n)}$ of a single
Polyakov loop is also an important variable. As we will discuss in Chap.~12, it
is an order parameter for the deconfinement transition in gluodynamics at
finite temperature.

\section{The static quark potential}

Having introduced the Wilson loop as an observable for the static quark
potential $V(r)$, we now discuss the general form of $V(r)$. As a first step we
calculate the potential in the limit of strong coupling $g$ (corresponding to
small $\beta$ -- compare (3.5)) and find that this gives rise to a linearly
rising term. Subsequently, we argue that, for small coupling $g$, we obtain the
$1/r$ potential familiar from electrodynamics. Thus, we find that the static
QCD potential can be parameterized by
\begin{equation}
  V(r) = A + \frac{B}{r} + \sigma r .
\end{equation}
Since the force between the quarks is the derivative of $V(r)$, the constant
$A$ is only an irrelevant normalization of the energy. The second term in
(3.62) is the Coulomb part of the potential with strength $B$. Finally, the
third contribution is a linearly rising term and the real constant $\sigma$ is
the so-called string tension. From QCD phenomenology one expects a value of
$\sigma \approx 900\ \mathrm{MeV}/\mathrm{fm}$. After providing the evidence
for the individual terms in (3.62), we discuss the physical implications of the
static QCD potential.

\subsection{Strong coupling expansion of the Wilson loop}

In order to demonstrate the presence of the linearly rising term in (3.62), we
calculate the vacuum expectation value of the Wilson loop in the limit of
strong coupling, i.e., large $g$ (small $\beta$). More explicitly we compute
\begin{equation}
  \ev{W_C} = \frac{1}{Z}\int D[U]\,
    \exp\!\left[-\frac{\beta}{3}\sum_P \Re\,\tr[1 - U_P]\right]
    \tr\!\left[\prod_{l\in C} U_l\right] .
\end{equation}
For this calculation we use a simplified notation: The sum runs over all
plaquettes $P$, where each plaquette is counted with only one of the two
possible orientations. The product over $l$ runs over all link variables
contained in the contour $C$ defining the Wilson loop. This expression can be
rewritten as
\begin{equation}
\begin{split}
  \ev{W_C} &= \frac{1}{Z'}\int D[U]\,
    \exp\!\left[\frac{\beta}{3}\sum_P \Re\,\tr[U_P]\right]
    \tr\!\left[\prod_{l\in C} U_l\right] \\
  &= \frac{1}{Z'}\int D[U]\,
    \exp\!\left[\frac{\beta}{6}\sum_P \left(\tr[U_P] + \tr[U_P^\dagger]\right)\right]
    \tr\!\left[\prod_{l\in C} U_l\right] .
\end{split}
\end{equation}
In the first step we separate the constant factor
$\exp\!\left(-\frac{\beta}{3}\sum_P \Re\,\tr[1]\right)$ from the Boltzmann
factor $\exp(-S)$. Exactly the same constant factor appears in the partition
function $Z$ and we cancel the two factors in the numerator and the
denominator. The partition function without this factor is denoted by $Z'$.
In the second step we use
\begin{equation}
  \Re\,\tr[U_P] = \frac{1}{2}\left(\tr[U_P] + \tr[U_P^\dagger]\right) .
\end{equation}
We stress that, according to (2.48), hermitian conjugation of the plaquette
variable $U_P$ is equivalent to inverting the orientation of the plaquette.
Thus in the second line of (3.64) we explicitly display both orientations of
the plaquette variables $U_P$, which leads to an extra factor $1/2$.

In the form of (3.64) we can now discuss the expansion of the Wilson loop
expectation value for strong coupling (small $\beta$). In particular we expand
the Boltzmann factor of (3.64) in $\beta$ using the Taylor expansion for the
exponential function,
\begin{equation}
  \begin{split}
    \exp\!\left[\frac{\beta}{6}\sum_P \left(\tr[U_P] + \tr[U_P^\dagger]\right)\right]
    &= \sum_{i,j=0}^{\infty}\frac{1}{i!\,j!}\left(\frac{\beta}{6}\right)^{i+j}\\
    &\qquad \left(\sum_P \tr[U_P]\right)^{i}
      \left(\sum_P \tr[U_P^\dagger]\right)^{j} .
  \end{split}
\end{equation}
Note that in this expansion we have separated the contributions from clockwise
oriented plaquettes $U_P^\dagger$ and counter-clockwise oriented plaquettes
$U_P$. This is important since for the leading term in the expansion only those
plaquettes oriented oppositely to the Wilson loop contribute.

For the normalization factor $Z'$ it is straightforward to determine the
leading contribution in the small-$\beta$ expansion. Already the first term with
$i = j = 0$ in (3.66) gives a nonvanishing contribution to the integral and we
obtain (using the normalization of the Haar measure)
\begin{equation}
  \begin{split}
    Z' &= \int D[U]\,\exp\!\left[\frac{\beta}{6}\sum_P
      \left(\tr[U_P] + \tr[U_P^\dagger]\right)\right]\\
    &= \int D[U]\,\left(1 + \order{\beta}\right)
    = 1 + \order{\beta^2} .
  \end{split}
\end{equation}
The expansion of the numerator of (3.64) is less straightforward. If only the
leading term in the expansion of the Boltzmann factor is kept, then the product
of link variables building up the observable, $\prod_{l\in C} U_l$, gives rise
to integrals of the type (3.22) which all vanish. Thus, in order to find the
leading nonvanishing term of the expansion we have to expand the Boltzmann
factor in small $\beta$. This brings down additional link variables from the
exponent and in this way we can saturate the integrals over the links to obtain
nonvanishing contributions of the type (3.24). In sub-leading terms (3.25) also
contributes.

If we consider the contour $C$ of the Wilson loop to be a $n_r \times n_t$
rectangle of links, then the minimal area $A_C$ spanned by this contour
contains $n_A = n_r n_t$ plaquettes (note that $n_r$, $n_t$, $n_A$ are positive
integers). The physical area $A_C$ is related to the extension of the Wilson
loop in physical units $a n_r$, $a n_t$ by $A_C = a^2 n_A = a n_r\, a n_t$.
Remembering (3.29) and Fig.~3.1 we find nonvanishing contributions only when
each link variable $U_\mu(n)$ in the loop is paired with its conjugate partner
$U_\mu(n)^\dagger$. Since we have plaquettes in our action, this must continue
until we have filled the contour $C$ with $n_A$ plaquettes obtained from the
expansion of the Boltzmann factor. We depict this contribution in Fig.~3.4.

Note that the plaquettes used for filling the contour have to have the opposite
orientation of the Wilson loop. Only in this way the contributions at each link
have the form of the integrand in (3.24). Since we need at least
$n_A = n_r n_t$ plaquettes from the exponent, the necessary term in the
expansion (3.66) of the exponential is of order $n_A$. Explicitly this leading
term reads (note that from the two orientations in (3.66) only the one opposite
to the Wilson loop contributes)
\begin{equation}
\begin{split}
  &\int D[U]\;\frac{1}{n_A!}\left(\frac{\beta}{6}\right)^{n_A}
    \left(\sum_P \tr[U_P^\dagger]\right)^{n_A}
    \tr\!\left[\prod_{l\in C} U_l\right] \\
  &\qquad = \left(\frac{\beta}{6}\right)^{n_A}\int D[U]\;
    \left(\sum_{P\in A_C} \tr[U_P^\dagger]\right)
    \tr\!\left[\prod_{l\in C} U_l\right] \\
  &\qquad = \left(\frac{\beta}{6}\right)^{n_A}\tr[\mathbb{1}]
    \left(\frac{1}{3}\right)^{n_A}
  = 3\,\exp\!\left[n_A \ln\!\left(\frac{\beta}{18}\right)\right] .
\end{split}
\end{equation}
In the first step of (3.68) we expand the $n_A$-th power over the sum of
plaquette variables with the correct orientation ($U_P^\dagger$). We keep only
the terms where each of the $n_A$ plaquettes $P$ inside the minimal area $A_C$
is occupied by the matching $U_P^\dagger$. There are exactly $n_A!$ such
products and thus the factor $1/n_A!$ is canceled. All other terms in the
expansion of the $n_A$-th power vanish since they give rise to integrals of the
type (3.22). In the second step we evaluate the nonvanishing term using the
integral (3.29) in the form depicted in Fig.~3.1. We glue together the
plaquette variables inside the Wilson loop first into rows and then these rows
into the full inner block of Fig.~3.4. Another gauge integral ties this block
to the oppositely oriented outer contour of the Wilson loop. One finds that
these steps give rise to the factor $(1/3)^{n_A}$ in (3.68). Combining (3.67)
and (3.68) we find
\begin{equation}
  \ev{W_C} = 3\,\exp\!\left[n_A \ln\!\left(\frac{\beta}{18}\right)\right]
    \left(1 + \order{\beta}\right)
  = 3\,\exp\!\left[n_r n_t \ln\!\left(\frac{\beta}{18}\right)\right]
    \left(1 + \order{\beta}\right) .
\end{equation}
According to (3.56) this expression has to be compared to the asymptotic form,
i.e., for large $t = a n_t$ we have
\begin{equation}
  \ev{W_C} \propto \exp\left(-a n_t V(r)\right) .
\end{equation}
Thus, we conclude that in the strong coupling limit (note that $r = a n_r$)

\begin{figure}[htbp]
  \centering
  \includegraphics[width=0.8\textwidth]{fig34.png}
  \caption*{Fig. 3.4. \emph{Leading contribution in the strong coupling (small
  $\beta$) expansion of the Wilson loop. The outer, counter-clockwise-oriented
  rectangle is the Wilson loop, the smaller, clockwise-oriented squares are the
  single plaquette terms from the action}.}
  \label{fig:34}
\end{figure}
\begin{equation}
  V(r) = \sigma r ,
\end{equation}
where the string tension $\sigma$ is given by the leading order expression
\begin{equation}
  \sigma = -\frac{1}{a^2}\ln\!\left(\frac{\beta}{18}\right)\left(1 + \order{\beta}\right) .
\end{equation}
We remark that it is relatively easy to produce higher corrections to the string
tension. However, since the strong coupling expansion does not play a central
role in modern lattice gauge theory (we will soon see that we are actually
interested in approaching weak coupling), we will not discuss the calculation
of higher terms and the convergence properties of this series.\footnote{The
original derivation by Wilson [15] has been rigorously proven in [16]. We refer
the reader to Creutz' book [1] for an elementary presentation or to [17] for a
more advanced calculation.} Instead, we stress once more that with a relatively
simple expansion we have extracted from the lattice formulation a linearly
rising potential. Below we discuss that such a term in the potential gives rise
to the important feature of confinement. However, before we come to the
discussion of the physical implications, let us first present the argument for
the presence of the Coulomb-type term in the parameterization (3.62) of the
static QCD potential.

\subsection{The Coulomb part of the static quark potential}

The presence of a Coulomb-type interaction can be readily seen from the
behavior of the gluon action in the continuum for small coupling constant $g$.
This argument is best presented in the normalization of the gauge field
introduced in (2.18) and (2.19), where the gauge field $A_\mu$ was rescaled by
a factor of $1/g$. In this form the Euclidean continuum gauge field action,
written as a sum over the color components, reads (compare (2.23) and (2.24))
\begin{equation}
  S_G[A] = \frac{1}{4}\sum_{i=1}^{8}\int \dd^4 x\,
    F^{(i)}_{\mu\nu}(x)\,F^{(i)}_{\mu\nu}(x) ,
\end{equation}
where the field strength tensor for individual color components is given by
\begin{equation}
  F^{(i)}_{\mu\nu}(x)
  = \partial_\mu A^{(i)}_{\nu}(x) - \partial_\nu A^{(i)}_{\mu}(x)
  - g\, f^{ijk} A^{(j)}_{\mu}(x)\,A^{(k)}_{\nu}(x) .
\end{equation}
From (3.74) it is obvious that the self-interaction terms that are specific for
QCD are multiplied by the coupling constant $g$. Thus, when sending this
coupling to zero, the field strength tensor (3.74) reduces to its abelian
counterpart, i.e., it has the form known from QED. Thus, the action (3.73)
turns into a sum over QED-type interactions for each color component. Since we
know that in QED the static potential is of the Coulomb-type,\footnote{See [18]
for an elementary derivation of this fact in the path integral formalism.} we
conclude that the $1/r$ term should be present in the parameterization (3.62)
of the static QCD potential.

\subsection{Physical implications of the static QCD potential}

Let us briefly sketch the physical implications of the static QCD potential, in
particular the role of the linear term. The linearly rising term in the
potential between a static quark--antiquark pair implies that the energy keeps
rising linearly as one tries to pull the two constituents apart. Thus the quark
and the antiquark are confined in a strongly bound meson state. Similarly, as
we show later, also a combination of three quarks is bound, forming a baryon.
The phenomenon that only color neutral combinations like hadrons are observable
objects is one possible definition of confinement.

The physical mechanism which leads to the linearly rising term is the formation
of a flux tube between the two sources. In QED, where we have no
self-interaction of the gauge field, the field lines between a source and a
sink spread out in space. In QCD the strong self-interaction of the gluons
prevents this behavior, and the field is squeezed into a narrow tube or string
producing the linear rise. Direct experimental evidence for the linearly rising
potential is seen when the mass of hadrons is plotted as a function of their
total spin and a linear behavior is found (see, e.g., the discussion in [19]).
Since for a linearly rising potential the energy rises linearly with the
angular momentum, this experimental finding confirms the linear term in (3.62).

So far our discussion is based only on the static potential obtained from pure
gluodynamics. Certainly also the quarks, which we have not yet included in our
discussion, play an important role. In the full theory with dynamical quarks,
processes of particle--antiparticle creation and annihilation become important.
In particular if the quark and antiquark are pulled sufficiently far apart, the
energy becomes large enough to create a quark--antiquark pair, which may
recombine with the two initial constituents to form two mesons. This phenomenon
is called string breaking and can be studied on the lattice.

To summarize, we find that lattice QCD is very friendly to confinement. This
property can easily be proven in the strong coupling limit, but as we shall see
in Chap.~4, it is also relatively simple to extract the QCD potential in a
numerical calculation at weaker coupling.

\section{Setting the scale with the static potential}

In the lattice formalism the action is given in units of $\hbar$ and all
observables are dimensionless. Only by relating them to physical quantities we
may introduce a scale parameter. An example of such a dimensionless quantity
calculated on the lattice is the product $aM$ of the lattice spacing $a$ and
some mass $M$, which may be determined from an exponential decay of a
correlation function like in (1.21). Identifying $M$ with a physical mass
allows us to determine the lattice constant $a$ in physical units.

In this section we discuss another method for setting the scale, which is based
on the static quark potential $V(r)$. This allows us to relate the lattice
spacing $a$ to the inverse gauge coupling $\beta$. We define a certain physical
distance $r_0$, the so-called Sommer parameter [20], which is a characteristic
length scale tied to the static potential. In physical units this distance is
$r_0 \approx 0.5$ Fermi (1 Fermi $= 1$ fm $= 10^{-15}$ m). This distance allows
us to determine the lattice spacing $a$ from $V(r)$ simply by counting the
number of lattice points between $r = 0$ and $r = r_0$. At the end of this
section we will be able to discuss the true continuum limit of pure lattice
gauge theory.

\subsection{Discussion of numerical data for the static potential}

Let us begin our presentation with a discussion of numerical data for the
static potential. At the moment we assume that the data have already been
calculated since we introduce the actual techniques for such a calculation only
in Chap.~4. We assume that we have computed the expectation value $\ev{W_C}$
for a planar Wilson loop of size $r \times t$, i.e., the contour $C$ is an
$r \times t$ rectangle. The spatial distance $r$ and the Euclidean time $t$ are
related to integer numbers $n$ and $n_t$ through the lattice spacing $a$,
\begin{equation}
  r = n a , \qquad t = n_t a .
\end{equation}
According to (3.56) the vacuum expectation value of the Wilson loop is
connected to the static potential via (we here omit the exponentially
suppressed corrections shown explicitly in (3.56))
\begin{equation}
  \ev{W_C} = C\,\exp\left(-t V(r)\right)
  = C\,\exp\left(-n_t a\,V(na)\right) .
\end{equation}
Since at this stage the lattice spacing $a$ is still unknown, all we can extract
from our numerical data for $\ev{W_C}$ are values for the product $aV(an)$ at
different $n$. In practice this is done via a two-parameter fit of the data for
different $n_t$ (but fixed $n$) according to (3.76). The fit parameters are $C$
and $aV(an)$. This procedure can be repeated for different values of $n$, and
the result is a set of numerical data for $aV(an)$ as a function of $n$.

In Fig.~3.5 we show the results from a numerical calculation of $aV(an)$ as a
function of $n$ for two different values of $\beta$. The actual numerical
results are represented by the symbols which we connect with dotted lines to
guide the eye. Since the data are from a numerical simulation, they come with
(small) statistical errors represented by horizontal bars. Note that in the two
plots different scales are used for the axes. The vertical scales are chosen
such that the shape of the static potential is similar in the two plots.

The plots nicely show the linearly rising behavior at larger values of $n$
which we have already deduced from our strong coupling calculation. For short
distances the data points also display a curvature which comes from the Coulomb
part of the potential. This is more pronounced in the right-hand side plot
which is for a larger value of $\beta$. In particular at small distances the
lattice Coulomb potential deviates from the continuum form. The lattice Coulomb
potential is related to the lattice-free boson propagator [23] and is often
used to correct for this deviation [24--26].

\begin{figure}[htbp]
  \centering
  \includegraphics[width=0.9\textwidth]{fig35.png}
  \caption*{Fig. 3.5. \emph{Numerical data for the static quark potential
  computed with the Wilson gauge action at two different couplings $\beta$. The
  symbols are the numerical data which we connect to guide the eye. The dashed
  vertical lines are drawn at a distance that corresponds to the Sommer
  parameter. The data are taken from [21, 22]}.}
  \label{fig:35}
\end{figure}

\subsection{The Sommer parameter and the lattice spacing}

We have already announced that we introduce the Sommer parameter as a way to
determine the lattice spacing $a$. Later in the book we discuss alternative
methods to determine $a$, e.g., by using the masses of hadrons. Although the
different methods provide values of $a$ that differ slightly at a given gauge
coupling, in the continuum limit the final physical results should agree.

The Sommer parameter $r_0$ is a certain distance related to the shape of the
static potential. Its physical value is $r_0 \approx 0.5$ fm. Although we
discuss the precise definition of the distance $r_0$ only below, we have
already marked its position in the two plots of Fig.~3.5 using dashed vertical
lines. Thus, the distance between $n = 0$ on the horizontal axis and the dashed
line corresponds to 0.5 fm. Note that the variable $n$ on the horizontal axis
is dimensionless -- it simply is the number of lattice spacings in the spatial
direction of the contour $C$. Thus, what is marked by the vertical dashed lines
is the Sommer scale in lattice units, i.e., $r_0/a$. We will discuss below how
to extract the ratio $r_0/a$ from the numerical data for $aV(an)$.

Once the Sommer scale $r_0/a$ is computed, it is easy to determine the lattice
spacing. When inspecting the two plots in Fig.~3.5, we find that the number of
lattice points between the origin ($n = 0$) and the vertical dashed line at
$r_0/a$ is different for the two values of $\beta$. In the left-hand side plot
of Fig.~3.5, corresponding to $\beta = 5.95$, we find about $n = 4.9$ lattice
points to the left of the dashed line. In the plot for $\beta = 6.20$ we find
approximately $n = 7.4$ lattice points left of the vertical line. Since in both
cases the distance between $n = 0$ and the vertical line corresponds to 0.5 fm,
we find for the lattice spacing $a$:
\begin{equation}
\begin{aligned}
  &a = 0.5\ \mathrm{fm}\,/\,4.9 \approx 0.102\ \mathrm{fm}
    && \text{for }\beta = 5.95 ,\\
  &a = 0.5\ \mathrm{fm}\,/\,7.4 \approx 0.067\ \mathrm{fm}
    && \text{for }\beta = 6.20 .
\end{aligned}
\end{equation}
We finally need to discuss how the Sommer scale $r_0$ is defined [20]. The
Sommer scale is not based directly on the potential $V(r)$, but instead on the
force $F(r) = dV(r)/dr$ between the two static quarks.\footnote{This definition
of $F$ differs by a minus sign from the usual definition.} For sufficiently
heavy quarks, quark--antiquark bound states can be described by an effective
nonrelativistic Schr\"odinger equation and the force $F(r)$ can be studied.
From comparing with experimental data for the $b\bar b$ and $c\bar c$ spectra
one finds that
\begin{equation}
  F(r_0)\,r_0^2 = 1.65 \qquad \text{corresponds to} \quad r_0 \approx 0.5\ \mathrm{fm} .
\end{equation}
Thus, we need to calculate from our numerical data for $V(r)$ the dimensionless
product $F(r)\,r^2$ and determine the value $r = r_0$ where this product
assumes the value 1.65. For an easy-to-digest presentation of this
determination we discuss the Sommer parameter directly for the potential as
parameterized in (3.62). We obtain for the force
\begin{equation}
  F(r) = \frac{d}{dr}V(r)
  = \frac{d}{dr}\left(A + \frac{B}{r} + \sigma r\right)
  = -\frac{B}{r^2} + \sigma .
\end{equation}
Thus, for the parameterized potential the condition (3.78) reads
\begin{equation}
  F(r_0)\,r_0^2 = -B + \sigma r_0^2 = 1.65 ,
\end{equation}
which implies $r_0 = \sqrt{(1.65 + B)/\sigma}$, or expressed in lattice units
\begin{equation}
  \frac{r_0}{a} = \sqrt{\frac{1.65 + B}{\sigma a^2}} .
\end{equation}
The ratio $r_0/a$ is exactly what we have used in Fig.~3.5 and the subsequent
determination of $a$. The numbers $B$ and $\sigma a^2$ can be determined from
the fit\footnote{We remark that such a fit already assumes that the potential
has -- at least locally -- the form (3.62). An alternative approach is to
directly determine the force $F$ by discretizing the derivative $d/dr$ on the
lattice (see [20] for details). This approach has the further advantage that one
does not need to fit $\sigma a^2$ and $B$, a fit which is dominated by large
values of $n$ where the signal to noise ratio of the data for $aV(an)$ becomes
poor.} of our data for $aV(an)$ to (this is equation (3.62) for $r = an$)
\begin{equation}
  aV(an) = Aa + \frac{B}{n} + \sigma a^2 n ,
\end{equation}
for different values of $n$. Thus, the determination of the lattice spacing can
be summarized as follows: Step 1: Determine $B$ and $\sigma a^2$ from the
numerical data for $aV(an)$. Step 2: Use (3.81) to calculate the dimensionless
number $X = r_0/a$. Step 3: The lattice spacing is then given by
$a = (0.5/X)\ \mathrm{fm}$.

In the discussion of Fig.~3.5 we have already addressed the fact that the
lattice spacing decreases when increasing $\beta$ (see also (3.77)). In [22]
the lattice spacing $a$ was determined for the Wilson action at several values
of $\beta$ and the dependence of $a$ on $\beta$ was parametrized for
$5.7 \leq \beta \leq 6.92$ as
\begin{equation}
  a = r_0\,\exp\!\left[-1.6804 - 1.7331(\beta - 6) + 0.7849(\beta - 6)^2
    - 0.4428(\beta - 6)^3\right] .
\end{equation}
The form of the parametrization in (3.83) is inspired by the renormalization
group, an important concept which we briefly address now.

\subsection{Renormalization group and the running coupling}

Coupling constants like the gauge coupling $g$ or a quark mass $m$ as they
enter the action functional are usually called bare parameters. These are not
directly observable ``physical'' numbers. Only by computing observables such as
hadron masses, the string tension, or the Sommer parameter, and by identifying
those with experimental values one can find out the values of the bare
parameters of the action in physical units.

Lattice actions may differ in various aspects. They may use different
discretizations of derivatives or the lattice grid, which is usually taken to
be hypercubic, may vary in its structure. However, when removing the lattice
cutoff, i.e., in the limit $a \to 0$, physical observables should agree with
the experimental value and become independent of $a$. In general this will
imply that the bare parameters have a nontrivial dependence on the cutoff $a$,
meaning that they are functions $g(a)$, $m(a)$, etc. As we send $a \to 0$ the
values of the bare parameters will have to be changed in order to keep physics
constant. This running of the bare parameters is addressed by the so-called
renormalization group. To simplify the discussion of this idea we consider pure
gauge theory where we have only one bare coupling, the gauge coupling $g(a)$.

Let $P(g(a), a)$ be a physical observable which in the limit $a \to 0$ obtains
its physical value $P_0$,
\begin{equation}
  \lim_{a\to 0} P(g(a), a) = P_0 .
\end{equation}
Callan and Symanzik, following early suggestions for QED by St\"uckelberg,
Peterman, Gell-Mann, and Low, formulated the requirement of constant physics in
a differential equation:
\begin{equation}
  \frac{dP(g, a)}{d\ln a} = 0 \qquad \text{or, equivalently} \qquad
  \left(\frac{\partial}{\partial\ln a} + \frac{\partial g}{\partial\ln a}
    \frac{\partial}{\partial g}\right) P(g, a) = 0 .
\end{equation}
The equation relates to a semi-group of scale-changing transformations, hence
the name renormalization group equation. (Actually, and more precisely, the
right-hand side of this equation is $\order{(a/\xi)^2 \ln(a/\xi)}$ for a
lattice system with correlation length $\xi$.) The coefficient function of the
second term is called the $\beta$-function (not to be mixed up with our inverse
gauge coupling),
\begin{equation}
  \beta(g) \equiv -\frac{\partial g}{\partial\ln a} ,
\end{equation}
and determines, up to an integration constant, how the coupling $g$ depends on
the cutoff $a$.

The $\beta$-function may be expanded in a power series around $g = 0$, with
coefficients determined by perturbation theory. For SU($N$) and $n_f$ massless
quarks the result reads
\begin{equation}
\begin{aligned}
  \beta(g) &= -\beta_0\,g^3 - \beta_1\,g^5 + \order{g^7} ,\\
  \beta_0 &= \frac{1}{(4\pi)^2}\left(\frac{11}{3}N - \frac{2}{3}n_f\right) ,\\
  \beta_1 &= \frac{1}{(4\pi)^4}\left(\frac{34}{3}N^2 - \frac{10}{3}N\,n_f
    - \frac{N^2 - 1}{N}\,n_f\right) .
\end{aligned}
\end{equation}
For pure gauge SU(3) we have to take $N = 3$ and $n_f = 0$. These first two
coefficients of the expansion are universal, independent of the regularization
scheme. In general, however, the $\beta$-function will depend on the details of
the regularization.

The differential Eq.~(3.86) with (3.87) can be solved, using separation of
variables, and one obtains
\begin{equation}
  a(g) = \frac{1}{\Lambda_L}\left(\beta_0 g^2\right)^{-\beta_1/(2\beta_0^2)}
    \exp\!\left(-\frac{1}{2\beta_0 g^2}\right)\left(1 + \order{g^2}\right) .
\end{equation}
The integration constant $\Lambda_L$ is used to set the scale by fixing the
value of $g$ at some $a$. Inverting the relation (3.88) one obtains the
coupling $g$ as a function of the scale $a$, the so-called running coupling,
\begin{equation}
  g(a)^{-2} = \beta_0 \ln\!\left(a^{-2}\Lambda_L^{-2}\right)
  + \frac{\beta_1}{\beta_0}\ln\!\left(\ln\!\left(a^{-2}\Lambda_L^{-2}\right)\right)
  + \order{1/\ln\!\left(a^2 \Lambda_L^2\right)} .
\end{equation}
Changing $a$ thus implies a corresponding change of $g$ such that physical
observables remain independent of the scale-fixing procedure. The value of
$\Lambda_L$ depends on the regularization scheme. Different continuum or lattice
actions have different values of $\Lambda_L$; their ratios may be exactly
related by a 1-loop perturbative calculation, however.

In fact, the validity of the perturbative expansion (3.87) has to be checked.
We find that for shrinking lattice spacing the running coupling also decreases
(for $n_f < 11N/2$). Vanishing lattice spacing corresponds to vanishing coupling
$g$ since $g = 0$ is indeed a zero of the $\beta$-function. This behavior is
called asymptotic freedom. For honesty we have to point out that there are
possible nonperturbative contributions of the type $\order{\exp(-1/g^2)}$
(nonvanishing for $g > 0$ but not contributing to the series expansion). It is
generally expected that these are not relevant for sufficiently small $g$. Also,
the order of performing the infinite volume limit and the continuum limit is an
important issue, cf. [27] and references therein.

\subsection{The true continuum limit}

At the end of Chap.~1, we already briefly mentioned how the true continuum
limit of a lattice field theory can be taken. However, at that point we did not
have the prerequisites to address this limit in a proper way and promised a
detailed discussion later. For the case of pure gauge theory, we can now provide
this discussion: Since we have shown in (3.83) and (3.88), respectively, that
the lattice spacing $a$ decreases with decreasing $g$ (increasing $\beta$), we
conclude that we simply have to study the limit
\begin{equation}
  \beta \to \infty
\end{equation}
to obtain the true continuum limit $a \to 0$. There are, however, certain
caveats to be considered in this procedure. If one performs the limit (3.90),
then the physical volume of the box in which we study QCD is proportional to
$a^4$ and thus shrinks to zero, unless we also increase the numbers of lattice
points in the spatial ($N$ points) and temporal ($N_T$ points) directions of
our lattice. In an ideal world one would first perform the so-called
thermodynamic limit
\begin{equation}
  N \to \infty , \qquad N_T \to \infty ,
\end{equation}
and only after that step the continuum limit (3.90) would be taken. However,
since in a numerical calculation this is not feasible, one is reduced to
calculating the physical observables for a few values of $\beta$, giving rise to
different values of $a$. The numbers of lattice points $N$, $N_T$ are always
chosen such that the physical extension
\begin{equation}
  L = a N , \qquad T = a N_T ,
\end{equation}
of the box remains fixed for the different values of $a$. Studying the
$a$-dependence of the results at fixed physical volume allows one to analyze
the dependence on the scale $a$ and to extrapolate the results to $a \to 0$.
The study of the $a$-dependence is often referred to as scaling analysis. The
extrapolation to $a = 0$ can then be repeated for different physical sizes $L$,
$T$ which in the end allows one to extrapolate the data to infinite physical
volume.

\section{Lattice gauge theory with other gauge groups}

Although this book is dedicated to lattice QCD, where the gauge group is SU(3),
we also briefly discuss lattice gauge theory with other gauge groups, in
particular SU($N$) and U(1). SU(2) is of interest since it is simpler than
SU(3) and is also a subgroup of SU(3), and it is therefore widely studied.
Important ideas, such as topological charge and instantons, which are solutions
of the classical equations of motion for the gauge field, can be formulated in
SU(2). The groups SU($N$) with $N \geq 4$ are of interest since in the
continuum the limit $N \to \infty$ is a limit appealing for analytic studies.
Thus, on the lattice one can analyze the $N$-dependence of observables. For
general $N$ the Wilson action reads
\begin{equation}
  S_G[U] = \frac{\beta}{N}\sum_{n\in\Lambda}\sum_{\mu<\nu}
    \Re\,\tr\!\left[1 - U_{\mu\nu}(n)\right] ,
\end{equation}
with $\beta = 2N/g^2$. The plaquette $U_{\mu\nu}(n)$ is the same ordered product
(2.48) as for SU(3). All observables which we have discussed so far, in
particular Wilson and Polyakov loops, can be taken over unchanged to general
$N$.

The gauge group U(1) is interesting since it corresponds to QED. As an abelian
group it is simpler to implement, and a numerical simulation of this case is a
fairly easy exercise that can produce interesting results on a PC within a few
hours. The main difference to the non-abelian groups is the fact that the
elements of U(1) are not matrices but simply complex phases, i.e., we can write
the link variables as $U_\mu(n) = \exp(\im A_\mu(n))$, where the $A_\mu(n)$ are
real. This implies that the commutator of two link variables vanishes. Hence,
the quadratic term in the field strength vanishes and one finds
$F_{\mu\nu}(x) = \partial_\mu A_\nu(x) - \partial_\nu A_\mu(x)$. No
self-interaction terms emerge. The Wilson action for U(1) is given by
\begin{equation}
  S_G[U] = \beta\sum_{n\in\Lambda}\sum_{\mu<\nu}
    \Re\!\left(1 - U_{\mu\nu}(n)\right) .
\end{equation}
No trace over color degrees of freedom is necessary and the plaquette is simply
a product of complex numbers,
$U_{\mu\nu}(n) = U_\mu(n)\,U_\nu(n+\hat\mu)\,U_\mu(n+\hat\nu)^*\,U_\nu(n)^*$
(the asterisk denotes complex conjugation). The coupling $\beta$ is related to
the coupling $e$ of QED by $\beta = 1/e^2$. The Polyakov loop and Wilson loop
observables can again be used, but no trace is needed.

An interesting observation can be made when applying the strong coupling
(small-$\beta$) expansion to U(1) lattice gauge theory. One still finds a
linearly rising term in the potential, indicating confinement. On the other hand
this theory is expected to describe QED, which does not show confinement. The
solution of this riddle is a phase transition in the U(1) lattice gauge theory.
At a critical value of $\beta_{\mathrm{crit}} \approx 1.01$ the theory changes
its behavior such that above $\beta_{\mathrm{crit}}$ the linearly rising term
of the potential vanishes and only the Coulomb potential term remains. For
SU($N$) lattice gauge theories it is generally believed that no such transition
occurs and the theory remains confining for all $\beta$. However, for an
alternative scenario see [28].

\section*{References}

\begin{enumerate}[label={[\arabic*]}]
  \item M.~Creutz: \emph{Quarks, Gluons and Lattices} (Cambridge University
        Press, Cambridge, New York 1983).
  \item H.~Georgi: \emph{Lie Algebras in Particle Physics}
        (Benjamin/Cummings, Reading, Massachusetts 1982).
  \item H.~F.~Jones: \emph{Groups, Representations and Physics} (Hilger,
        Bristol 1990).
  \item M.~Hamermesh: \emph{Group Theory and Its Application to Physical
        Problems} (Addison-Wesley, Reading, Massachusetts 1964).
  \item R.~Gilmore: \emph{Lie Groups, Lie Algebras, and Some of Their
        Applications}, Vol.~64 (John Wiley \& Sons, New York 1974).
  \item C.~Itzykson and J.-B.~Zuber: \emph{Quantum Field Theory} (McGraw-Hill,
        New York 1985).
  \item M.~E.~Peskin and D.~V.~Schroeder: \emph{An Introduction to Quantum
        Field Theory} (Addison-Wesley, Reading, Massachusetts 1995).
  \item S.~Weinberg: \emph{The Quantum Theory of Fields}, Vol.~1 and 2
        (Cambridge University Press, Cambridge, New York 1996).
  \item W.~H.~Press, S.~A.~Teukolsky, W.~T.~Vetterling, and B.~P.~Flannery:
        \emph{Numerical Recipes in C}, 2nd ed. (Cambridge University Press,
        Cambridge, New York 1999).
  \item L.~Giusti et al.: Int. J. Mod. Phys. A 16, 3487 (2001).
  \item S.~Elitzur: Phys. Rev. D 12, 3978 (1975).
  \item T.~DeGrand and C.~DeTar: \emph{Lattice Methods for Quantum
        Chromodynamics} (World Scientific, Singapore 2006).
  \item A.~M.~Polyakov: Phys. Lett. B 59, 82 (1975).
  \item C.~Borgs and E.~Seiler: Commun. Math. Phys. 91, 329 (1983).
  \item K.~G.~Wilson: Phys. Rev. D 10, 2445 (1974).
  \item K.~Osterwalder and E.~Seiler: Ann. Phys. 110, 440 (1978).
  \item J.~M.~Drouffe and J.~B.~Zuber: Phys. Rep. 102, 1 (1983).
  \item G.~Roepstorff: \emph{Path Integral Approach to Quantum Physics}
        (Springer, Berlin, Heidelberg, New York 1996).
  \item D.~H.~Perkins: \emph{Introduction to High Energy Physics}, 4th ed.
        (Cambridge University Press, Cambridge, New York 1999).
  \item R.~Sommer: Nucl. Phys. B 411, 839 (1994).
  \item M.~Guagnelli, R.~Sommer, and H.~Wittig: Nucl. Phys. B 535, 389 (1998).
  \item S.~Necco and R.~Sommer: Nucl. Phys. B 622, 328 (2002).
  \item C.~B.~Lang and C.~Rebbi: Phys. Lett. B 115, 137 (1982).
  \item C.~Michael: Nucl. Phys. B 259, 58 (1985).
  \item R.~G.~Edwards, U.~M.~Heller, and T.~R.~Klassen: Nucl. Phys. B 517, 377
        (1998).
  \item C.~R.~Allton et al.: Phys. Rev. D 65, 054502 (2002).
  \item E.~Seiler: in \emph{Applications of RG Methods in Mathematical
        Sciences} (RIMS, Kyoto University, Kyoto 2003).
  \item A.~Patrascioiu and E.~Seiler: Phys. Rev. Lett. 74, 1920 (1995).
\end{enumerate}
